\subsection{Diffuse measures; atomic measures}

DEFINITION 5. --- A measure \(\theta\) on T is said to be diffuse if \(|\theta|(\{t\}) = 0\) for every \(t \in \mathrm{T}\) .

Example. --- Lebesgue measure on \(\mathbf{R}\) is diffuse (Ch. IV, §1, No.~3, Remark 1).

To say that \(\theta\) is a diffuse measure on T amounts to saying that every set with finite complement carries \(|\theta|\) , or again that \(|\theta|\) is alien to every point measure. The diffuse measures therefore form a band in \(\mathcal{M}(\mathrm{T})\) (Ch. II, §1, No.~5, Th. 1).

Recall (Ch. III, §1, No.~3) that a complex measure \(\rho\) on T is said to be atomic if it is of the form \(\sum_{t\in \mathrm{T}}\alpha (t)\varepsilon_t\) , where \(\alpha\) is a complex function on T such that \(\sum_{t\in \mathrm{K}}|\alpha (t)| < +\infty\) for every compact subset K of T, which expresses that the family \(\left(\alpha (t)\varepsilon_t\right)_{t\in \mathrm{T}}\) is summable (§2, No.~1, Remark 2). It then follows from the remark following Cor. 3 of Th. 2 of No.~5 that \(|\rho| = \sum_{t\in \mathrm{T}}|\alpha (t)|\varepsilon_t\) . The function \(\alpha\) that occurs in these formulas is uniquely determined, because \(\alpha (t) = \rho (\{t\})\) . An atomic measure and a diffuse measure are alien to each other.

PROPOSITION 15. --- Every complex measure \(\sigma\) on T may be written in one and only one way in the form \(\rho + \theta\) , where \(\rho\) is an atomic measure and \(\theta\) is a diffuse measure; one then has \(|\sigma| = |\rho| + |\theta|\) .

The uniqueness of the decomposition is obvious since, \(\rho\) being atomic and \(\theta\) diffuse, necessarily \(\rho = \sum_{t\in \mathrm{T}}\rho (\{t\})\varepsilon_t = \sum_{t\in \mathrm{T}}\sigma (\{t\})\varepsilon_t\) and \(\theta = \sigma -\rho\) . To establish existence, it suffices to observe that \(\sum_{t\in \mathrm{K}}|\sigma (\{t\})| \leqslant |\sigma |(\mathrm{K}) < +\infty\) for every compact set \(\mathbf{K}\) , so that one can set \(\sum_{t\in \mathrm{T}}\sigma (\{t\})\varepsilon_t = \rho\) . The measure \(\sigma -\rho\) is obviously diffuse, and the relation \(|\sigma | = |\rho | + |\sigma -\rho |\) follows at once from Cor. 3 of Prop. 13 of No.~7.

One observes that this proof shows that if \(\sigma\) is carried by a set \(M\) and if \(|\sigma| (\{t\}) > 0\) for every \(t \in M\) , then \(\sigma\) is atomic.

